\documentclass[12pt,a4paper]{article}
\usepackage[utf8]{inputenc}
\usepackage[bulgarian]{babel}
\usepackage{amsmath}
\usepackage{amsfonts}
\usepackage{amssymb}
\usepackage{graphicx}
\usepackage{geometry}
\usepackage{fancyhdr}
\usepackage{listings}
\usepackage{xcolor}
\usepackage{tikz}
\usepackage{array}
\usepackage{booktabs}
\usepackage{multirow}

\geometry{margin=2.5cm}
\pagestyle{fancy}
\fancyhf{}
\rhead{Тема 13}
\lhead{Компютърни мрежи и протоколи}
\rfoot{Страница \thepage}

\lstset{
    basicstyle=\ttfamily\small,
    breaklines=true,
    frame=single,
    columns=fullflexible
}

\begin{document}

\title{\textbf{Тема 13: Компютърни мрежи и протоколи}\\ \large{Изпит}}
\author{}
\date{}
\maketitle

\tableofcontents
\newpage

\section{OSI модел и TCP/IP}

\subsection{OSI модел — седемте слоя}

\textbf{Идея:} Разделяне на сложната мрежова комуникация на независими слоеве, всеки с ясна роля.

\begin{center}
\begin{tabular}{|l|l|}
\hline
\textbf{7. Application (Приложен)} & HTTP, FTP, DNS \\
\hline
\textbf{6. Presentation (Презентационен)} & TLS, кодиране \\
\hline
\textbf{5. Session (Сесиен)} & управление \\
\hline
\textbf{4. Transport (Транспортен)} & TCP, UDP \\
\hline
\textbf{3. Network (Мрежов)} & IP, маршрути \\
\hline
\textbf{2. Data Link (Канален)} & Ethernet, Wi-Fi \\
\hline
\textbf{1. Physical (Физически)} & кабели, сигнали \\
\hline
\end{tabular}
\end{center}

\vspace{0.5cm}

\textbf{Ключови слоеве:}

\begin{center}
\begin{tabular}{|l|p{5cm}|l|l|}
\hline
\textbf{Слой} & \textbf{Роля} & \textbf{PDU} & \textbf{Устройства} \\
\hline
7-Application & Интерфейс към приложенията & Данни & — \\
\hline
4-Transport & Доставка между процеси, надеждност & Сегмент/Дейтаграма & — \\
\hline
3-Network & Маршрутизация между мрежи, IP адреси & Пакет & Рутери \\
\hline
2-Data Link & Локална доставка, MAC адреси, грешки & Рамка & Суичове \\
\hline
1-Physical & Физически сигнали & Битове & Хъбове, кабели \\
\hline
\end{tabular}
\end{center}

\vspace{0.5cm}

\textbf{Енкапсулация} — всеки слой добавя свой header:

\begin{lstlisting}
Application:  [DATA]
Transport:    [TCP header][DATA]
Network:      [IP header][TCP header][DATA]
Data Link:    [Eth header][IP][TCP][DATA][CRC]
\end{lstlisting}

\subsection{Съпоставка OSI $\leftrightarrow$ TCP/IP}

\textbf{TCP/IP} е практическият модел с 4 слоя:

\begin{center}
\begin{tabular}{|l|c|l|}
\hline
\textbf{OSI (7 слоя)} & & \textbf{TCP/IP (4 слоя)} \\
\hline
7. Application & & \\
6. Presentation & $=$ & Application \\
5. Session & & \\
\hline
4. Transport & $=$ & Transport (TCP/UDP) \\
\hline
3. Network & $=$ & Internet (IP) \\
\hline
2. Data Link & & Network Access \\
1. Physical & $=$ & (Ethernet, Wi-Fi) \\
\hline
\end{tabular}
\end{center}

\vspace{0.5cm}

\textbf{Разлики:}
\begin{itemize}
    \item OSI — теоретичен, референтен модел (1984)
    \item TCP/IP — практически, развит с интернет (1970-те)
    \item OSI има повече слоеве, TCP/IP ги комбинира
\end{itemize}

\newpage
\section{Разпределена маршрутизация}

\textbf{Маршрутизация} = намиране на път от източник до получател през множество рутери.

\textbf{Защо разпределена?} Интернет е огромен — няма централен контролер. Всеки рутер взима собствени решения.

\subsection{Distance Vector (дистантен вектор)}

\textbf{Принцип:} Всеки рутер знае само \textbf{разстоянията до съседите} си и обменя своята таблица с тях.

\textbf{Алгоритъм:} Белман-Форд

\begin{lstlisting}
Всеки рутер:
1. Започва с разстояние 0 до себе си, ∞ до другите
2. Периодично изпраща ЦЕЛИЯ си вектор на съседите
3. Когато получи вектор от съсед N:
   D(x) = min(D(x), cost(N) + D_N(x))
4. Ако има промяна → обновява таблицата
\end{lstlisting}

\textbf{Пример:}

\begin{lstlisting}
A ── 1 ── B ── 1 ── C

Първоначално:
A: {A:0, B:1}
B: {A:1, B:0, C:1}
C: {B:1, C:0}

След обмен:
A: {A:0, B:1, C:2}  ← научава за C през B
\end{lstlisting}

\textbf{Предимства:}
\begin{itemize}
    \item[$\checkmark$] Прост
    \item[$\checkmark$] Малко памет
\end{itemize}

\textbf{Недостатъци:}
\begin{itemize}
    \item[$\times$] Бавна сходимост
    \item[$\times$] \textbf{Count-to-infinity problem}
\end{itemize}

\subsubsection{Count-to-infinity проблем}

\begin{lstlisting}
A ── 1 ── B ── 1 ── C

Връзката B-C се счупва:
1. B вижда C:∞
2. Но A казва "C:2 през B"
3. B мисли "мога през A за 3"
4. A чува "C:3 от B", мисли "мога за 4"
5. ∞ цикъл...
\end{lstlisting}

\textbf{Решения:}
\begin{itemize}
    \item Split horizon — не казвай на съсед маршрут, научен от него
    \item Poison reverse — кажи "$\infty$ през теб"
    \item Дефиниране на $\infty$ като малко число (напр. 16)
\end{itemize}

\textbf{Протоколи:} RIP, IGRP

\subsection{Link State (следене на състоянието на линията)}

\textbf{Принцип:} Всеки рутер изгражда \textbf{пълна карта} на топологията и самостоятелно изчислява най-краткия път.

\textbf{Алгоритъм:}

\begin{lstlisting}
Всеки рутер:
1. Открива съседите си (Hello пакети)
2. Измерва цената на връзките
3. Създава LSA (Link State Advertisement):
   "Аз съм R, имам линии към A(цена 3), B(цена 5)"
4. Излъчва (flooding) LSA до ВСИЧКИ рутери
5. Получава LSA от всички → строи пълна карта
6. Пуска алгоритъм на Дейкстра → намира най-краткия път
\end{lstlisting}

\textbf{Визуализация:}

\begin{lstlisting}
    A ─3─ B
    │     │
    4     2
    │     │
    C ─1─ D

Всеки рутер получава:
- A: "линии към B(3), C(4)"
- B: "линии към A(3), D(2)"
- C: "линии към A(4), D(1)"
- D: "линии към B(2), C(1)"

→ Всеки има пълната карта → пуска Дейкстра
\end{lstlisting}

\textbf{Дейкстра} (накратко):
\begin{enumerate}
    \item Започни от себе си (разстояние 0)
    \item На всяка стъпка: вземи най-близкия непосетен възел
    \item Обнови разстоянията на съседите му
    \item Повтаряй докато всички са посетени
\end{enumerate}

\textbf{Предимства:}
\begin{itemize}
    \item[$\checkmark$] Бърза сходимост
    \item[$\checkmark$] Няма count-to-infinity
    \item[$\checkmark$] По-добра мащабируемост
\end{itemize}

\textbf{Недостатъци:}
\begin{itemize}
    \item[$\times$] По-сложен
    \item[$\times$] Повече памет и CPU
\end{itemize}

\textbf{Протоколи:} OSPF, IS-IS

\subsection{Сравнение}

\begin{center}
\begin{tabular}{|l|l|l|}
\hline
& \textbf{Distance Vector} & \textbf{Link State} \\
\hline
\textbf{Знание} & Само разстояния до съседи & Пълна топология \\
\hline
\textbf{Обмен} & Цяла таблица със съседи & LSA към всички \\
\hline
\textbf{Алгоритъм} & Белман-Форд & Дейкстра \\
\hline
\textbf{Сходимост} & Бавна & Бърза \\
\hline
\textbf{Памет} & Малко & Повече \\
\hline
\textbf{Count-to-$\infty$} & Да (проблем) & Не \\
\hline
\end{tabular}
\end{center}

\newpage
\section{IPv4 адресация}

\subsection{Класова адресация (Classful)}

Старият подход (1981) — адресите се делят на класове:

\begin{lstlisting}
Клас A:  0....... |  Network (8 бита) | Host (24 бита)
         0-127     |  128 мрежи         | 16M хостове

Клас B:  10...... |  Network (16)      | Host (16)
         128-191   |  16K мрежи         | 65K хостове

Клас C:  110..... |  Network (24)      | Host (8)
         192-223   |  2M мрежи          | 254 хостове

Клас D:  1110.... |  Multicast
Клас E:  1111.... |  Резервиран
\end{lstlisting}

\textbf{Пример:}
\begin{itemize}
    \item \texttt{10.5.7.9} — клас A, мрежа \texttt{10.0.0.0}, хост \texttt{.5.7.9}
    \item \texttt{172.16.5.7} — клас B, мрежа \texttt{172.16.0.0}, хост \texttt{.5.7}
    \item \texttt{192.168.1.7} — клас C, мрежа \texttt{192.168.1.0}, хост \texttt{.7}
\end{itemize}

\textbf{Проблем:} Класове B са твърде големи, C — твърде малки $\rightarrow$ масово хабене на адреси.

\subsection{Безкласова адресация (CIDR)}

\textbf{CIDR} = Classless Inter-Domain Routing (1993)

\textbf{Идея:} Произволна дължина на мрежовата част, записвана с \textbf{префикс}.

\begin{lstlisting}
192.168.10.0/24    ← 24 бита мрежа, 8 бита хост = 256 адреса
10.0.0.0/8         ← 8 бита мрежа, 24 бита хост = 16M адреса
172.16.0.0/12      ← 12 бита мрежа, 20 бита хост
203.0.113.16/28    ← 28 бита мрежа, 4 бита хост = 16 адреса
\end{lstlisting}

\textbf{Subnet mask} — алтернативен запис:
\begin{itemize}
    \item /24 = \texttt{255.255.255.0}
    \item /28 = \texttt{255.255.255.240}
\end{itemize}

\subsubsection{Изчисляване на мрежа}

\textbf{Задача:} \texttt{192.168.10.130/26} — каква е мрежата, broadcast, диапазон?

\begin{lstlisting}
/26 → 26 бита мрежа, 6 бита хост
Маска: 255.255.255.192 (последният октет: 11000000 = 192)

192.168.10.130 в бинар:
11000000.10101000.00001010.10000010
                          ||||||
                          └─────┘ 6 бита хост

AND с маска:
11000000.10101000.00001010.10000000 = 192.168.10.128 ← мрежа

Broadcast (всички 1 в host):
11000000.10101000.00001010.10111111 = 192.168.10.191

Използваеми: 192.168.10.129 - 192.168.10.190 (62 адреса)
\end{lstlisting}

\subsubsection{Subnetting (подмрежи)}

\textbf{Задача:} \texttt{192.168.1.0/24} $\rightarrow$ раздели на 4 подмрежи

\begin{lstlisting}
4 подмрежи = 2² → нужни 2 бита → нов префикс /26

Резултат:
192.168.1.0/26    (0-63)
192.168.1.64/26   (64-127)
192.168.1.128/26  (128-191)
192.168.1.192/26  (192-255)
\end{lstlisting}

\textbf{Специални адреси:}
\begin{itemize}
    \item \texttt{127.0.0.0/8} — loopback
    \item \texttt{10.0.0.0/8}, \texttt{172.16.0.0/12}, \texttt{192.168.0.0/16} — частни мрежи (RFC 1918)
\end{itemize}

\newpage
\section{IPv6 — основни характеристики}

\subsection{Защо IPv6?}

IPv4 има само \textbf{4.3 милиарда} адреси $\rightarrow$ изчерпани към 2011 г.

IPv6 има $\mathbf{2^{128} \approx 3.4 \times 10^{38}}$ адреси $\rightarrow$ практически неизчерпаемо.

\subsection{Формат на адреса}

\textbf{128 бита}, записвани като 8 групи по 16 бита (hex):

\begin{lstlisting}
2001:0db8:85a3:0000:0000:8a2e:0370:7334

Съкращения:
1. Махни водещи нули: 2001:db8:85a3:0:0:8a2e:370:7334
2. Замени поредица от нули с ::
   2001:db8:85a3::8a2e:370:7334
   (само ВЕДНЪЖ в адреса!)

::1       — loopback (като 127.0.0.1)
::        — unspecified
\end{lstlisting}

\subsection{Ключови подобрения}

\begin{center}
\begin{tabular}{|l|l|l|}
\hline
\textbf{Характеристика} & \textbf{IPv4} & \textbf{IPv6} \\
\hline
Адресно пространство & 32 бита (4.3B) & 128 бита (340 ундецилиона) \\
\hline
Header & променлив, 20-60 байта & фиксиран, 40 байта \\
\hline
Фрагментация & от рутери & само от източника \\
\hline
Checksum в IP & да & не (на по-горни слоеве) \\
\hline
ARP & да & не (заменен с NDP) \\
\hline
Auto-configuration & DHCP & SLAAC (без сървър) \\
\hline
IPsec & опционален & задължителен (на хартия) \\
\hline
Broadcast & да & не (multicast/anycast) \\
\hline
NAT & широко използван & ненужен \\
\hline
\end{tabular}
\end{center}

\subsection{Типове адреси}

\begin{lstlisting}
Unicast (един получател):
├─ Global unicast (2000::/3)     — публични, рутируеми
├─ Link-local (fe80::/10)         — само в локалния линк
└─ Unique local (fc00::/7)        — частни

Multicast (ff00::/8)               — група получатели
Anycast                            — най-близкия от група
\end{lstlisting}

\subsection{Преход IPv4 $\rightarrow$ IPv6}

Не е "big bang" — съжителстват:

\begin{enumerate}
    \item \textbf{Dual-stack} — устройството поддържа и двата
    \item \textbf{Тунелиране} — IPv6 пакети в IPv4 (6to4, Teredo)
    \item \textbf{Превод} — NAT64, DNS64
\end{enumerate}

\newpage
\section{TCP — трикратно договаряне}

\subsection{Защо TCP?}

IP е \textbf{ненадежден} — пакети могат да се загубят, дублират, разменят.

\textbf{TCP} осигурява:
\begin{itemize}
    \item[$\checkmark$] Надеждна доставка (препредаване)
    \item[$\checkmark$] Подреждане (sequence numbers)
    \item[$\checkmark$] Контрол на потока
    \item[$\checkmark$] Контрол на задръстванията
    \item[$\checkmark$] Connection-oriented комуникация
\end{itemize}

\subsection{Three-way handshake (3-етапно ръкостискане)}

\textbf{Цел:} Установяване на връзка и обмен на начални sequence numbers (ISN).

\begin{lstlisting}
Клиент (C)                          Сървър (S)
   │                                    │
   │  (1) SYN                           │
   │  seq=x ─────────────────────────►  │  "Искам връзка, започвам с x"
   │                                    │
   │  (2) SYN+ACK                       │
   │  ◄────────── seq=y, ack=x+1        │  "ОК, започвам с y, очаквам x+1"
   │                                    │
   │  (3) ACK                           │
   │  ack=y+1 ───────────────────────►  │  "Потвърждавам, очаквам y+1"
   │                                    │
   │    ESTABLISHED                     │
   │  ◄──────── данни ──────────────►   │
\end{lstlisting}

\subsubsection{Стъпка 1: SYN (Synchronize)}

Клиентът:
\begin{itemize}
    \item Избира \textbf{случаен ISN} $x$
    \item Изпраща: \texttt{SYN=1, seq=x}
    \item Състояние: \textbf{SYN\_SENT}
\end{itemize}

\subsubsection{Стъпка 2: SYN-ACK}

Сървърът:
\begin{itemize}
    \item Избира \textbf{свой ISN} $y$
    \item Изпраща: \texttt{SYN=1, ACK=1, seq=y, ack=x+1}
    \item Състояние: \textbf{SYN\_RECEIVED}
    \item Казва: "Чух те и също искам връзка"
\end{itemize}

\subsubsection{Стъпка 3: ACK}

Клиентът:
\begin{itemize}
    \item Изпраща: \texttt{ACK=1, seq=x+1, ack=y+1}
    \item Състояние: \textbf{ESTABLISHED}
    \item Може да носи вече данни (piggybacking)
\end{itemize}

Сървърът получава $\rightarrow$ \textbf{ESTABLISHED}

\subsection{Защо точно ТРИ стъпки?}

\textbf{Защо не две?}

Защото връзката е \textbf{бидирекционална} — всяка посока има свой sequence number и трябва да бъде потвърдена.

\begin{quote}
\textit{Аналогия:}\\
"Чуваш ли ме?"\\
\hspace*{1cm}"Да, чувам. А ти чуваш ли МЕН?"\\
\hspace*{2cm}"Да, чувам."
\end{quote}

След стъпка 2: клиентът знае, че двата канала работят\\
След стъпка 3: сървърът знае, че двата канала работят

\textbf{Защита от стари пакети:} Три стъпки предотвратяват установяване на невалидна сесия от стар SYN пакет.

\subsection{Затваряне — 4-way handshake}

\begin{lstlisting}
C ── FIN ──────►  S    "Приключих"
C ◄── ACK ────── S    "Чух те"
C ◄── FIN ────── S    "И аз приключих"
C ── ACK ──────► S    "Чух те"
\end{lstlisting}

(понякога стъпки 2 и 3 се комбинират)

\subsection{SYN Flood атака}

Атакуващ изпраща много SYN-и с фалшиви IP $\rightarrow$ сървърът алокира ресурси $\rightarrow$ опашката се препълва $\rightarrow$ легитимни клиенти не могат да се свържат.

\textbf{Защита:} SYN cookies — не алокира памет до финалния ACK.

\newpage
\section{DNS — резолвинг на имена}

\subsection{Какво е DNS?}

\textbf{DNS} = Domain Name System — превръща имена (\texttt{www.example.com}) в IP адреси.

\textbf{Защо?} Хората помнят имена, компютрите — IP адреси.

\textbf{Структура:} Разпределена, йерархична, кеширана база данни.

\subsection{Йерархия}

\begin{lstlisting}
                       . (root)
                       │
        ┌──────┬───────┼───────┬──────┐
       com    org     net     bg    ...    ← TLD
        │                       │
     google                uni-sofia
        │                       │
       www                      fmi
                                 │
                                www
\end{lstlisting}

\textbf{Четене на име:} \texttt{www.fmi.uni-sofia.bg.}
\begin{itemize}
    \item \texttt{.} — root (обикновено скрит)
    \item \texttt{bg} — Top-Level Domain (TLD)
    \item \texttt{uni-sofia} — Second-Level Domain
    \item \texttt{fmi} — субдомейн
    \item \texttt{www} — хост
\end{itemize}

\subsection{Процес на резолвинг}

\textbf{Пример:} Искаме \texttt{www.example.com}

\begin{lstlisting}
   Браузър
      │
      ▼
   OS (stub resolver)
      │
      ▼
   Recursive resolver (8.8.8.8)
      │
      ├─► (1) Root сървър
      │      "За .com → питай TLD сървър X"
      │
      ├─► (2) TLD .com сървър
      │      "За example.com → питай Authoritative Y"
      │
      ├─► (3) Authoritative сървър (example.com)
      │      "www.example.com = 93.184.216.34"
      │
      ▼
   Връща IP на клиента и кешира
\end{lstlisting}

\textbf{Видове заявки:}
\begin{itemize}
    \item \textbf{Recursive} (рекурсивна) — клиент $\rightarrow$ resolver, resolver прави цялата работа
    \item \textbf{Iterative} (итеративна) — resolver $\rightarrow$ root/TLD/auth, всеки дава референция
\end{itemize}

\subsection{DNS записи}

\begin{center}
\begin{tabular}{|l|l|p{5cm}|}
\hline
\textbf{Тип} & \textbf{Значение} & \textbf{Пример} \\
\hline
A & IPv4 адрес & \texttt{www.example.com → 93.184.216.34} \\
\hline
AAAA & IPv6 адрес & \texttt{www.example.com → 2606:2800:220:1::1} \\
\hline
CNAME & Псевдоним & \texttt{www → example.com} \\
\hline
MX & Mail сървър & \texttt{example.com → mail.example.com} \\
\hline
NS & Name server & \texttt{example.com → ns1.example.com} \\
\hline
PTR & Обратен (IP→име) & \texttt{34.216.184.93.in-addr.arpa → www} \\
\hline
TXT & Текст (SPF, DKIM) & — \\
\hline
\end{tabular}
\end{center}

\textbf{Запомни:} \textbf{A} за IPv4, \textbf{AAAA} за IPv6 (4×A).

\subsection{Резолвинг по IPv4 и IPv6}

Един домейн може да има \textbf{и A, и AAAA} запис:

\begin{lstlisting}
Заявка за IPv4:
www.example.com A? → 93.184.216.34

Заявка за IPv6:
www.example.com AAAA? → 2606:2800:220:1::1
\end{lstlisting}

\textbf{Happy Eyeballs (RFC 8305):}\\
Модерните браузъри \textbf{паралелно}:
\begin{enumerate}
    \item Питат за A и AAAA
    \item Опитват и двата IP адреса едновременно
    \item Използват първия, който отговори успешно
\end{enumerate}

\textbf{Обратен резолвинг (Reverse DNS):}
\begin{itemize}
    \item IPv4: \texttt{34.216.184.93.in-addr.arpa} $\rightarrow$ PTR
    \item IPv6: всяка hex цифра обърната + \texttt{.ip6.arpa}
\end{itemize}

\subsection{Кеширане и TTL}

\textbf{TTL} (Time To Live) — за колко секунди може да се кешира записът.

Кеширане на нива:
\begin{itemize}
    \item Браузър
    \item Операционна система
    \item Recursive resolver
    \item Междинни сървъри
\end{itemize}

\textbf{Малък TTL} (300 сек) — бърза реакция при промяна, повече заявки\\
\textbf{Голям TTL} (86400 сек) — по-малко заявки, по-бавна промяна

\subsection{Транспорт}

\begin{itemize}
    \item \textbf{UDP, порт 53} — по подразбиране (бързо)
    \item \textbf{TCP, порт 53} — за големи отговори ($>512$ байта)
    \item \textbf{DoT} (DNS over TLS) — порт 853
    \item \textbf{DoH} (DNS over HTTPS) — порт 443
\end{itemize}

\subsection{Сигурност — DNSSEC}

\textbf{Проблем:} DNS отговорите могат да се фалшифицират (spoofing, cache poisoning).

\textbf{DNSSEC} добавя \textbf{криптографски подписи} към записите $\rightarrow$ resolver може да провери автентичността.

\textbf{Внимание:} DNSSEC \textbf{не криптира}, само подписва.

\newpage
\section{Обобщение}

\subsection{OSI vs TCP/IP}
\begin{itemize}
    \item OSI — 7 слоя (теоретичен модел)
    \item TCP/IP — 4 слоя (практически)
    \item Енкапсулация — всеки слой добавя header
\end{itemize}

\subsection{Маршрутизация}
\begin{itemize}
    \item \textbf{Distance Vector} — знае само разстояния, Белман-Форд, има count-to-$\infty$
    \item \textbf{Link State} — има пълна карта, Дейкстра, бърза сходимост
\end{itemize}

\subsection{IPv4 адресация}
\begin{itemize}
    \item \textbf{Класова} — A/B/C класове (остаряло)
    \item \textbf{CIDR} — произволен префикс (/24, /28...), subnet mask
    \item Subnetting — разделяне на мрежи
\end{itemize}

\subsection{IPv6}
\begin{itemize}
    \item 128 бита (340 ундецилиона адреси)
    \item По-прост header, автоконфигурация, няма NAT
    \item A запис за IPv4, AAAA за IPv6
\end{itemize}

\subsection{TCP Three-way handshake}
\begin{itemize}
    \item SYN $\rightarrow$ SYN-ACK $\rightarrow$ ACK
    \item Обмен на sequence numbers
    \item Бидирекционална синхронизация
\end{itemize}

\subsection{DNS}
\begin{itemize}
    \item Йерархична разпределена база (root $\rightarrow$ TLD $\rightarrow$ authoritative)
    \item Recursive resolver прави итеративни заявки
    \item A (IPv4) и AAAA (IPv6) записи
    \item Кеширане с TTL
\end{itemize}

\end{document}
